\subsection{The $f_3$-stress readout conjecture under wobble-broken translation weights}
\label{sec:q6-f3-stress-wobble-broken-translation-readout-cross-layer}

\paragraph{Record.}
This chapter records a finite conjectural interface between the $B^{*}_{Q6}$ synonymous residual coordinate and translation-layer measurements. The object is not a statement about protein function, mature abundance, localization, folded structure, organismal fitness, or a universal biological law. It is a restricted claim about whether an $f_3$-paired synonymous codon imbalance remains visible after amino-acid projection when the readout is confined to translation-adjacent rows: tAI, stAI, ribosome occupancy, or translation efficiency. The carrier is the finite coding-sequence-by-organism residual coordinate surface, restricted to the $f_3$-stress coordinate and to matched translation-readout rows. Its internal data are a relation, a spectrum, and a trigger condition: the relation compares the residual coordinate with a readout row; the spectrum records which coordinate directions survive controlled ablation; the trigger says that the $f_3$ row can become dynamical only where context tRNA supply or wobble weights distinguish a codon $c$ from its third-position partner $c^{f_3}$.

\paragraph{Finite codon boundary.}
The finite $f_3$ saturation row fixes the codon-level support before any biological interpretation is attempted. The mediated set $M=\operatorname{Med}(R)\cup\{\mathrm{ACA}\}$ has $20$ listed codons,
$$
\begin{aligned}
M=\{&\mathrm{AAA},\mathrm{AAG},\mathrm{ACA},\mathrm{ACG},\mathrm{AGA},\mathrm{AGG},\mathrm{AUA},\mathrm{AUG},\mathrm{CUA},\mathrm{CUC},\\
&\mathrm{CUG},\mathrm{CUU},\mathrm{UAA},\mathrm{UAG},\mathrm{UCA},\mathrm{UCG},\mathrm{UGA},\mathrm{UGG},\mathrm{UUA},\mathrm{UUG}\}.
\end{aligned}
$$
The corresponding saturated $f_3$ image of $R$ has $18$ listed codons,
$$
\begin{aligned}
\operatorname{Sat}_{f_3}(R)=\{&\mathrm{AAA},\mathrm{AAG},\mathrm{AGA},\mathrm{AGG},\mathrm{AUA},\mathrm{AUG},\mathrm{CUA},\mathrm{CUC},\mathrm{CUG},\\
&\mathrm{CUU},\mathrm{UAA},\mathrm{UAG},\mathrm{UCA},\mathrm{UCG},\mathrm{UGA},\mathrm{UGG},\mathrm{UUA},\mathrm{UUG}\}.
\end{aligned}
$$
The reported $f_3$ boundary list is empty. Thus the finite combinatorial row does not by itself provide a positive translation readout; it supplies only the residual coordinate on which a translation row may be tested.

\paragraph{Conjectural survival statistic.}
For each organism and each eligible translation readout $T_j$, let $S_{f3\text{-}stress,T_j}$ denote the controlled survival entry of the $f_3$-stress coordinate after amino-acid composition, organism identity, GC3, coding length, $M$-density, and mRNA controls are included. The intended support condition is not mere codon-topology alignment and not a modeled adaptation score in isolation. Support requires at least one strictly positive controlled entry, written schematically as
$$
\begin{aligned}
S_{f3\text{-}stress,T_j}>0,
\end{aligned}
$$
with the simpler codon baselines removed and the readout row still attached to its source measurement. The conjecture is broken if the powered panel remains non-positive or rank-degenerate, including the case where an all-$1.0$ rank matrix supplies no survival contrast. In that null case, the correct conclusion is absence of a detected translation-layer distinction under the stated controls, not irrelevance of the $Q6$ coordinate in every biological setting.

\paragraph{Mechanistic boundary from tRNA and wobble rows.}
The strongest available mechanism-adjacent finite row is the tRNA-poverty assay on the $R$ sense codons. It was computed for $40$ of $46$ requested organisms, skipped $6$, used $1000$ permutations, and applied $\alpha=0.05$. Its recorded call says that $R$ sense codons are systematically tRNA-poor under GtRNAdb tGCN combined with dos Reis wobble weights. This supports a tRNA-supply interpretation of $R$ avoidance, but it still does not observe ribosome motion or protein output directly. In the present conjecture it functions as a break condition for symmetry: if the contextual weight $w_c$ equals $w_{c^{f_3}}$, the $f_3$ pair has no translation-weight contrast at that context; if $w_c\ne w_{c^{f_3}}$, the coordinate is eligible for a translation-readout test.

\paragraph{Powered residual spectrum.}
The cross-organism residual spectrum supplies finite scale and orientation. In the per-organism axis-strength computation, $11$ of $16$ requested organisms were computed, with $5$ folds, at least $500$ proteins per organism, and $41$ synonymous-design columns. The captured spectra were led by $d_{\mathrm{resid4}}$ on average and then by $d_{\mathrm{tRNA}}$; the dominant-count row was $d_{\mathrm{resid4}}=11$. Domain means were explicitly descriptive because the panel was imbalanced, with counts $n=1$ archaeon, $n=2$ bacteria, and $n=8$ eukaryotes. The largest reported GC3 rank effect was the usage direction with $\rho=0.136$. These numbers keep the conjecture finite: the row is powered enough to define a residual spectrum, but the domain imbalance prevents a claim of broad biological universality.

\paragraph{Axis interpretation.}
Several powered rows constrain how much may be inferred from the spectrum. The optimal-direction tRNA-adaptation computation used $11$ of $16$ requested organisms, $5$ folds, at least $500$ proteins per organism, and $41$ synonymous-design columns. The non-tRNA residual optimality row concluded that the dominant residual axis aligns most with $f_3$ translation selection under the same computed-organism and design-column scale. The two-axis closure row concluded that a third axis is present, with the reason that $d_{\perp}$ has systematically positive held-out $R^2$ and mutually same-signed directions across organisms. The third-axis identity row identified $d_{\perp}$ provisionally with within-family usage frequency, with mean cosine $0.128$ and one-sided sign-test $p=0.000488$. The mechanism row then reported selection-not-GC, not wobble-clustered, and that CSC or positional identity still needs a further outside measurement. Together these rows justify a restricted residual-axis reading, not a completed causal mechanism.

\paragraph{Fourth-axis guard.}
The finite fourth-axis rows are useful precisely because they block premature promotion. The three-axis sufficiency assay, again at $11$ computed organisms out of $16$, $5$ folds, at least $500$ proteins, and $41$ synonymous-design columns, concluded that a fourth mechanism is present. The mRNA-stability assay concluded that the fourth axis is not mRNA stability in the tested yeast half-life alignment, because $d_{\mathrm{opt\_halflife}}$ had low alignment with the abundance-associated $d_{\mathrm{resid4}}$ under the fixed gates. The elongation-dwell row for \textit{Saccharomyces cerevisiae} was marked partial. Therefore the present chapter cannot replace translation-efficiency measurement with mRNA half-life, and it cannot identify elongation dwell as settled from the partial row alone.

\paragraph{Contacts and non-promotion rule.}
The admissible contacts are organism-level codon usage, tRNA-Leu copy-number spectra, tAI or stAI weights, ribosome profiling translation efficiency, and matched mRNA abundance. Codon-usage tables measure sequence composition. tRNA gene-copy rows summarize genome annotations. tAI and stAI weights are modeled translation-efficiency annotations. Ribosome profiling supplies per-CDS ribosome occupancy or translation efficiency, typically as an RPF-to-mRNA ratio, and matched RNA abundance supplies the control needed to separate translation-layer signal from transcript abundance. None of these contacts, singly or jointly, licenses a conclusion about folding, physical admissibility, mature protein abundance, function, phenotype, or a global biological law unless a distinct measurement at that layer is attached. The bounded readback for this record is only the controlled survival entry $S_{f3\text{-}stress,T_j}$ for tAI, stAI, ribosome occupancy, or translation efficiency.

\paragraph{Falsification criterion.}
The conjecture is falsified in its intended finite form when every eligible $B^{*}_{Q6}$ translation pair remains non-positive after the full control set, or when the available ranks collapse so that no contrast remains after amino-acid projection and organism matching. It is supported only when at least one translation readout row remains positive under the controls and under simpler-codon-baseline ablation. Even then, the support is local to the measured translation row and to the organisms in the powered panel; it is not a claim about downstream structure, function, localization, stability, phenotype, mechanism completeness, or universality.
